\section{Global support control and stabilization}\label{sec:support}
The main difficulty is that the nonfixed support $s$ is not bounded as $n$ grows. Fixed-support calculations alone do not justify summing the asymptotic expansion. We first obtain a uniform positive majorant for every support, then prove sharper behavior at each fixed support.

\subsection{A uniform cycle index majorant}
Substitute $p_1=x$ and $p_\ell=x^\ell t^\ell$ for $\ell\geq2$. The homogeneous cycle-index polynomial becomes
\begin{equation}\label{eq:Hjt}
\mathcal H_j(x,t)=h_j(x,x^2t^2,x^3t^3,\ldots)
 =x^j[z^j]\frac{\exp((1-t)z)}{1-tz}.
\end{equation}
Although the last expression contains $1-t$, its coefficients in $x,t$ are nonnegative by the original cycle-index definition. Set
\[
F(x,t)=(1-x)^{-1}G(x,t),\qquad
G(x,t)=\prod_{j\geq2}(1-\mathcal H_j(x,t))^{-1},
\qquad B_{n,s}=[x^nt^s]F(x,t).
\]
Thus $B_{n,s}=\sum_{\rho\vdash s,\,\min\rho\geq2}\alpha_{1^{n-s}\rho}$.
At $t_0=1/4$,
\begin{equation}\label{eq:C0}
\sum_{j\geq2}\mathcal H_j(1,t_0)=\frac43e^{3/4}-2<1,
\qquad
C_0=\frac1{3-\frac43e^{3/4}}<6.
\end{equation}
For nonnegative $v_j$ with sum less than $1$, the finite-product inequality
$\prod_j(1-v_j)\geq1-\sum_jv_j$ and monotone convergence imply
$G(1,t_0)\leq C_0$. Multiplication by $(1-x)^{-1}$ takes partial sums of the nonnegative $x$ coefficients of $G$, so
\begin{equation}\label{eq:Bns}
B_{n,s}\leq C_0\,4^s\qquad(n,s\geq0).
\end{equation}
This estimate is uniform even when $s$ is comparable with $n$.

\subsection{Root and involution bounds}
\begin{lemma}\label{lem:bounds}
For $k\geq1$, $n\geq1$, and $0\leq s\leq n$,
\begin{equation}\label{eq:ratbound}
\sqrt k\leq\frac{I_k}{I_{k-1}}\leq1+\sqrt k,
\qquad
\frac{I_{n-s}}{I_n}\leq\fall ns^{-1/2}\leq(e/n)^{s/2}.
\end{equation}
If $\rho$ has no $1$-parts and $|\rho|=s\geq1$, then
\begin{equation}\label{eq:rootbound}
r(\rho)\leq2^{s/2}s^{s/4}.
\end{equation}
\end{lemma}
\begin{proof}
Writing $R_k=I_k/I_{k-1}$, the recurrence gives $R_k=1+(k-1)/R_{k-1}$. Induction proves the first pair of bounds. For the lower one, with $u=\sqrt{k-1}$, use
\[
1+\frac{u^2}{1+u}=u+\frac1{1+u}\geq\sqrt{u^2+1};
\]
squaring leaves the nonnegative difference $u^2/(1+u)^2$. The upper one follows from $R_{k-1}\geq\sqrt{k-1}$. The initial case is $R_1=1$.

Multiplying the lower ratio bounds proves the first bound on $I_{n-s}/I_n$. The mean of the largest $s$ values of $\log k$, $1\leq k\leq n$, is at least the mean of all $n$ values. Consequently
$\fall ns\geq(n!)^{s/n}\geq(n/e)^s$, proving the second bound.

The ratio upper bound also gives $I_m\leq2^m m^{m/2}$ for $m\geq1$. The factor in \eqref{eq:rootcycles} at cycle length $\ell\geq2$ is at most $\ell^{m_\ell/2}I_{m_\ell}$. Hence
\[
r(\rho)\leq2^{\sum m_\ell}\prod_{\ell\geq2}(\ell m_\ell)^{m_\ell/2}
\leq2^{s/2}s^{s/4},
\]
because $\sum m_\ell\leq s/2$ and $\ell m_\ell\leq s$.
\end{proof}

Combining \eqref{eq:Bns}, \eqref{eq:rootfactor}, and Lemma~\ref{lem:bounds}, the normalized contribution at support $s\geq1$ is at most
\begin{align}
\frac{I_{n-s}}{I_n}
 \sum_{\substack{\rho\vdash s\\\min\rho\geq2}}
       \alpha_{1^{n-s}\rho}r(\rho)
&\leq C_0 4^s 2^{s/2}s^{s/4}(e/n)^{s/2}\notag\\
&\leq C_0\bigl(4\sqrt{2e}\,n^{-1/4}\bigr)^s.\label{eq:globaltail}
\end{align}
For all sufficiently large $n$ the base of this geometric bound is less than $1$. Thus for every fixed $T$ the total contribution from $s\geq T$ is
\begin{equation}\label{eq:tailT}
O_T(I_n n^{-T/4}).
\end{equation}
The estimate is deliberately rough. Its role is uniform control of all large supports, not an attractive numerical onset.

\subsection{Exponential stabilization at fixed support}
Keep separate marked variables, substituting $p_1=x$ and $p_\ell=x^\ell u_\ell$ for $\ell\geq2$. Define
\[
G(x,u)=\prod_{j\geq2}(1-h_j(x,x^2u_2,x^3u_3,\ldots))^{-1},
\quad
Q_s(x)=\sum_{\substack{\rho\vdash s\\\min\rho\geq2}}
                r(\rho)[u_\rho]G(x,u).
\]
The root-weighted cycle-index coefficient at support $s$, \emph{before} multiplication by $I_{n-s}$, is
\begin{equation}\label{eq:Qcoefficient}
[x^n]\frac{Q_s(x)}{1-x}.
\end{equation}

\begin{lemma}\label{lem:stabilization}
For each fixed $s$ and each $1<R<\sqrt2$, the function $Q_s$ is analytic in a neighborhood of $|x|\leq R$, and
\begin{equation}\label{eq:stabilize}
[x^n]\frac{Q_s(x)}{1-x}=CH_s+O_{s,R}(R^{-n}).
\end{equation}
Individually, for every fixed type $\rho$ with no $1$-parts,
\begin{equation}\label{eq:alphastabilize}
\alpha_{1^{n-|\rho|}\rho}=Cq_\rho+O_{\rho,R}(R^{-n}).
\end{equation}
\end{lemma}
\begin{proof}
At $u=0$ the denominators are $1-x^j/j!$. For $|x|\leq R<\sqrt2$, one has $R^j/j!<1$ for every $j\geq2$, with supremum strictly less than $1$. Mark all $u_\ell$ by $t^\ell$. For sufficiently small positive $t$, the positive majorants $\mathcal H_j(R,t)$ stay uniformly below $1$ and have summable total. To see this, \eqref{eq:Hjt} gives a finite summed series whenever $Rt<1$; continuity controls a finite initial set of factors as $t\downarrow0$, and its summable tail is uniformly small. Only finitely many marked variables can enter a fixed weight. The infinite product is therefore locally uniformly convergent and analytic on a slightly larger $x$ disk and a small polydisc in those variables. Cauchy coefficient extraction proves analyticity of $Q_s$ and of each individual marked coefficient.

The coefficient in \eqref{eq:Qcoefficient} is the partial sum of the Taylor coefficients of $Q_s$. The omitted tail is $O_{s,R}(R^{-n})$, so that partial sum tends to $Q_s(1)$. At $x=1$,
\[
h_j(1,u_2,u_3,\ldots)=\frac1{j!}
 +\sum_{\substack{\rho\text{ has parts }\geq2\\0<|\rho|\leq j}}
                  \frac{u_\rho}{(j-|\rho|)!z_\rho}.
\]
Factoring $1-1/j!$ from the denominator gives $G(1,u)=CP(u)$, and hence $Q_s(1)=CH_s$. The same argument without the root-weighted sum proves \eqref{eq:alphastabilize}.
\end{proof}

\subsection{Proof of the main theorem}
Fix $S\geq0$ and choose $T=2S+2$. At each retained support $s\leq S$, Lemma~\ref{lem:stabilization} gives
$I_{n-s}(CH_s+O(R^{-n}))$. The omitted but bounded supports $S<s<T$ contribute $O_S(I_n n^{-(S+1)/2})$, by the same lemma and \eqref{eq:ratbound}. The entire remaining range $s\geq T$ is bounded by \eqref{eq:tailT}, with precisely the same order. Exponentially small retained errors can be absorbed. This proves Theorem~\ref{thm:main}.

The second cutoff $T$ is important. Applying the crude global bound directly at $S+1$ would yield only $n^{-(S+1)/4}$, which is weaker than the natural fixed-support remainder.
